% ===================================================================
%  QUADRO N.º 8 — A Colheita. — A Caça. — A Vindima. — A Pesca.
%
%  Ceci n'est PAS une transcription : c'est une traduction, faite en
%  2026, en portugais d'aujourd'hui. Voir texto/pt/10-quadro-01.tex
%  pour la consigne, qui vaut pour les seize fichiers.
% ===================================================================

%%K t08-noto-1 noto
\VUnotes{143.2mm}{%
\textit{(*) Porque esta palavra não é precisa: trata-se de uma colheita
de linho, de uma vindima, de uma apanha de maçãs? Talvez fosse melhor
usar «} \VUgras{mesono} \textit{», proposto em} Mondo \textit{XIV, p.
242. Mas até agora essa raiz nova não foi adotada pela Academia e
espera a discussão segundo as regras habituais do movimento idista.}%
}

%%K t08-apar-1 apar
\VUsaut{5.0mm}
\VUtitre{15.0pt}{60}{QUADRO N.º 8}
\VUsaut{3.0mm}
\VUcentre{13.2pt}{90}{\textit{Primeira cena.}}
\VUsaut{3.0mm}
\VUpk{10.2pt}{40}{A Colheita (*).}
\VUsaut{4.0mm}
\VUpk{10.2pt}{40}{Os Aspetos do Campo.}
\VUsaut{3.0mm}

%%K t08-c1-01-1 p
1. --- Com o \VUgras{verão}, o campo mudou de aspeto mais uma vez. A
semente confiada ao sulco germinou; uma plantinha verde saiu da terra e
deu a sua espiga. O grão formou-se, depois amadureceu, e já não resta
senão colher.

%%K t08-c1-02-1 p
2. --- As \VUgras{flores do campo} estão abertas. \VUgras{Centáureas}
\textsuperscript{(1)}, \VUgras{papoilas} \textsuperscript{(2)} e \VUgras{dedaleiras}
\textsuperscript{(3)} põem no tom uniforme dos trigais o alegre contraste das
suas frescas \VUgras{corolas}.

%%K t08-c1-03-1 p
3. --- As árvores de fruto cobriram-se de folhas; o fruto amadureceu. A
pequena \VUgras{cerejeira} \textsuperscript{(4)} está carregada de belas
\VUgras{cerejas} \textsuperscript{(5)}, que as crianças colhem e comem com
deleite.

%%K t08-c1-04-1 p
4. --- O rebanho pode já passar todo o dia ao ar livre.

%%K t08-c1-04-2 p
Os \VUgras{cordeiros} \textsuperscript{(6)} cresceram e tornaram-se belas
\VUgras{ovelhas} \textsuperscript{(7)} e belos \VUgras{carneiros}
\textsuperscript{(8)} de espesso velo. Pastam tranquilamente a \VUgras{erva} do
\VUgras{pasto} \textsuperscript{(9)}, salpicado de \VUgras{margaridas}.

%%K t08-c1-04-3 p
Em breve começarão a tosquiar estes animais para ter a sua \VUgras{lã},
com a qual se fabrica o pano da nossa roupa.

%%K t08-tit-2 sub
\VUsaut{5.0mm}
\VUpk{10.2pt}{40}{O Pastor.}
\VUsaut{3.4mm}

%%K t08-c1-05-1 p
5. --- O \VUgras{pastor} \textsuperscript{(10)} não parece fazer-lhes muito caso.
Só se preocupa com a trovoada que passa no horizonte, e o
\VUgras{zagal} \textsuperscript{(13)}, que leva o \VUgras{cantil}
\textsuperscript{(14)} aos lábios, parece ainda mais despreocupado.

%%K t08-c1-06-1 p
6. --- O calor é sufocante; porque o \VUgras{cão} \textsuperscript{(15)} vem também
refrescar-se; lambe a água fresca do \VUgras{ribeiro} \textsuperscript{(16)} e
assusta uma pobre \VUgras{rã} \textsuperscript{(17)} que estava acaçapada na erva
da margem. Esta salta para a água clara onde florescem os
\VUgras{nenúfares} \textsuperscript{(18)}.

%%K t08-c1-07-1 p
7. --- Uma \VUgras{alvéola} \textsuperscript{(19)} banha-se junto à pequena
\VUgras{nascente} \textsuperscript{(20)} que brota entre duas rochas e se esconde
sob o \VUgras{silvado} \textsuperscript{(21)} e as \VUgras{moitas de
pilriteiro} \textsuperscript{(22)}.

%%K t08-tit-3 sub
\VUsaut{5.0mm}
\VUpk{10.2pt}{40}{A Colheita.}
\VUsaut{3.4mm}

%%K t08-c1-08-1 p
8. --- Para as crianças do campo, a ceifa do trigo é uma época muito
divertida. Dão alegres \VUgras{cambalhotas} \textsuperscript{(24)} sobre os
\VUgras{montes de palha} \textsuperscript{(23)}, ajudam os \VUgras{ceifeiros}
\textsuperscript{(26)} e, enquanto estes ceifam o \VUgras{trigal}
\textsuperscript{(27)} com as suas \VUgras{gadanhas} \textsuperscript{(25)}, bem afiadas
na \VUgras{pedra de amolar} \textsuperscript{(30)}, elas recolhem o trigo e
atam-no em \VUgras{gavelas} \textsuperscript{(31)} ou em \VUgras{feixes}
\textsuperscript{(32)}.

%%K t08-c1-09-1 p
9. --- Às vezes permite-se aos mais velhos cortar com uma \VUgras{foice}
\textsuperscript{(12)} os \VUgras{caules dourados} \textsuperscript{(28)} que trazem as
pesadas \VUgras{espigas} \textsuperscript{(29)} cheias de grão; e os pequenos,
enquanto guardam o \VUgras{gado} \textsuperscript{(33)} que pasta no \VUgras{prado}
\textsuperscript{(34)} à sombra da grande \VUgras{tília} \textsuperscript{(35)}, caçam as
belas \VUgras{borboletas} com a sua \VUgras{rede} \textsuperscript{(36)}.

%%K t08-c1-09-2 p
Alguns até sobem ao \VUgras{carro} \textsuperscript{(37)} para colocar as gavelas
que lhes passam com um \VUgras{forcado} \textsuperscript{(38)}.

%%K t08-c1-10-1 p
10. --- Por toda a \VUgras{planície} \textsuperscript{(39)}, onde levantam voo as
\VUgras{cotovias} \textsuperscript{(41)}, reina a animação. Todos estão ocupados
com a ceifa. Mas, enquanto os pobres continuam a usar as ferramentas
antigas --- \VUgras{gadanhas, foices} e \VUgras{manguais}
\textsuperscript{(40)} ---, os proprietários ricos mandam ceifar o seu trigo ou
o seu \VUgras{centeio} \textsuperscript{(43)} com uma \VUgras{ceifeira mecânica}
\textsuperscript{(42)}. Depois, sem necessidade de levar as gavelas ao celeiro,
fazem-se ali mesmo grandes \VUgras{medas} \textsuperscript{(44)}.

%%K t08-c1-11-1 p
11. --- Então traz-se uma \VUgras{debulhadora} \textsuperscript{(47)} que uma
\VUgras{locomóvel} \textsuperscript{(45)} move por meio de uma \VUgras{correia de
transmissão} \textsuperscript{(46)}, e assim o grão fica separado da
\VUgras{palha} sem sair do campo em que foi semeado.

%%K t08-tit-4 sub
\VUsaut{5.0mm}
\VUpk{10.2pt}{40}{A Trovoada.}
\VUsaut{3.4mm}

%%K t08-c1-12-1 p
12. --- Muitas vezes rebentam violentas \VUgras{trovoadas}
\textsuperscript{(53)} na época da ceifa. Primeiro ouvem-se ao longe os
estrondos do \VUgras{trovão}; levanta-se um forte \VUgras{vento do oeste}
\textsuperscript{(54)} que verga com gemidos as árvores mais grossas do
\VUgras{bosque} \textsuperscript{(49)}. Negros \VUgras{nuvarrões}
\textsuperscript{(56)} assaltam o céu; atravessam-nos os ziguezagues cegantes do
\VUgras{relâmpago} \textsuperscript{(55)}. A \VUgras{chuva} e às vezes o
\VUgras{granizo} começam a cair, esmagando a colheita pronta a ser
ceifada.

%%K t08-c1-13-1 p
13. --- Muitas vezes o raio incendeia algum edifício. Assim, no ano
passado o telhado do \VUgras{moinho de vento} \textsuperscript{(50)} ficou
reduzido a cinzas e duas das suas \VUgras{velas} \textsuperscript{(51)} ficaram
destroçadas.

%%K t08-c1-14-1 p
14. --- Ao cabo de umas horas volta a calma. Um brilhante \VUgras{arco-íris}
\textsuperscript{(52)} aparece no horizonte e a natureza retoma a sua vida,
interrompida um instante. Os camponeses, que se tinham refugiado nas
\VUgras{cabanas} \textsuperscript{(48)}, voltam ao trabalho e esforçam-se
valentemente por reparar os danos que acabam de sofrer.

%%K t08-tit-5 sub
\VUsaut{6.0mm}
\VUcentre{13.2pt}{90}{\textit{Segunda cena.}}
\VUsaut{3.6mm}

%%K t08-tit-6 sub
\VUpk{10.2pt}{40}{A Caça. --- A Vindima.}
\VUsaut{1.4mm}
\VUpk{10.2pt}{40}{A Pesca.}
\VUsaut{4.0mm}

%%K t08-c2-01-1 p
1. --- Por finais de \VUgras{agosto} ou por meados de \VUgras{setembro},
conforme as regiões, abre-se a caça. Mal os primeiros raios do sol
fizeram sair as \VUgras{lagartixas} \textsuperscript{(2)} do seu buraco, já o
\VUgras{caçador} \textsuperscript{(5)} se põe a caminho com os seus \VUgras{cães}
\textsuperscript{(4)}.

%%K t08-c2-02-1 p
2. --- Pôs o seu sólido \VUgras{fato de caça} \textsuperscript{(9)} de pano grosso
e umas \VUgras{polainas} de couro com as quais poderá andar pela erva
molhada onde se escondem os \VUgras{sapos} \textsuperscript{(1)}; pegou na sua
\VUgras{espingarda} \textsuperscript{(6)} e no seu \VUgras{surrão}
\textsuperscript{(10)}, e já partiu.

%%K t08-c2-03-1 p
3. --- O ardor da perseguição leva-o a uma vinha onde a vindima está em
pleno auge. Os filhos do vinhateiro, que costumam guardar as cabras no
caminho, hoje deixam-nas pastar à vontade e, depois de atarem a uma
estaca um velho \VUgras{bode} \textsuperscript{(3)} que não deixaria de aproveitar
a sua liberdade para fugir, dão-se também a sua parte de prazer. Um tira
um cacho de uvas de um \VUgras{cesto de vindima} \textsuperscript{(12)} e
diverte-se a fazer correr o \VUgras{sumo} \textsuperscript{(14)} para um
\VUgras{copo} \textsuperscript{(13)} para fazer mosto (vinho não fermentado). O
outro coroa-se com uma \VUgras{grinalda de folhas} \textsuperscript{(15)} que
acaba de trançar. Junto deles, uns \VUgras{pardais} \textsuperscript{(16)}
atrevidos chegam até ao próprio lagar a debicar as uvas.

%%K t08-c2-04-1 p
4. --- Enquanto as crianças se divertem, os homens trabalham. Os
\VUgras{vindimadores} \textsuperscript{(18)} levam a uva à adega em
\VUgras{cestos de vindima} \textsuperscript{(17)}; um \VUgras{tanoeiro}
\textsuperscript{(19)} arranja as \VUgras{pipas} \textsuperscript{(20)}, verifica se as
\VUgras{aduelas} \textsuperscript{(22)} e os \VUgras{fundos} \textsuperscript{(21)} estão em
bom estado e substitui os \VUgras{arcos} \textsuperscript{(23)} partidos. Foi
também ele que fez as grandes \VUgras{dornas} \textsuperscript{(25)} nas quais se
deita a \VUgras{uva} \textsuperscript{(26)} para que fermente.

%%K t08-c2-05-1 p
5. --- Quando a fermentação terminou, tira-se o vinho e põe-se o resto no
\VUgras{lagar} \textsuperscript{(28)} e, apertando pouco a pouco o grande
\VUgras{fuso} \textsuperscript{(29)}, espreme-se o sumo, que corre para um
\VUgras{tonelete} \textsuperscript{(32)}, e obtém-se ainda mais \VUgras{vinho}
\textsuperscript{(31)}.

%%K t08-tit-8 sub
\VUsaut{6.0mm}
\VUpk{10.2pt}{40}{Cerveja e Sidra.}
\VUsaut{3.4mm}

%%K t08-c2-06-1 p
6. --- Pela parede da \VUgras{adega} \textsuperscript{(27)} trepa um ramo de
\VUgras{lúpulo} \textsuperscript{(30)}. Aqui não é mais do que uma planta de
adorno, mas nalguns países, onde a videira é muito rara, o lúpulo
cultiva-se para fabricar a \VUgras{cerveja}.

%%K t08-c2-07-1 p
7. --- A pequena \VUgras{macieira} \textsuperscript{(33)} está carregada de maçãs
que, quando estiverem maduras, darão uma \VUgras{sidra} excelente. Na sua
\VUgras{gaiola} \textsuperscript{(34)}, o \VUgras{canário} \textsuperscript{(35)} e o
\VUgras{pintassilgo} \textsuperscript{(36)} olham com inveja os pardais que
folgam em liberdade.

%%K t08-c2-08-1 p
8. --- Até onde a vista alcança, na encosta das colinas, alinham-se as
\VUgras{estacas} \textsuperscript{(40)} que sustentam as magníficas \VUgras{cepas}
\textsuperscript{(39)}, carregadas de pesados \VUgras{cachos}.

%%K t08-c2-08-2 p
Durante todo o ano, o \VUgras{vinhateiro} \textsuperscript{(42)} trabalhou para
tratar da sua \VUgras{vinha} \textsuperscript{(38)}, e agora vê-se recompensado do
seu esforço com uma bela colheita. As \VUgras{vindimadeiras}
\textsuperscript{(41)} enchem as dornas postas no carro que puxam as
\VUgras{mulas} \textsuperscript{(43)}, e todos estão alegres.

%%K t08-tit-9 sub
\VUsaut{5.0mm}
\VUpk{10.2pt}{40}{A Pesca.}
\VUsaut{3.4mm}

%%K t08-c2-09-1 p
9. --- O mesmo acontece com os \VUgras{pescadores de linha}
\textsuperscript{(44)}, que desde manhã vieram instalar-se à beira de água com
as suas \VUgras{canas de pesca} \textsuperscript{(45)}.

%%K t08-c2-09-2 p
O \VUgras{peixe} \textsuperscript{(47)} morde o anzol assim que deitam a sua
\VUgras{linha} \textsuperscript{(46)} à água, e mal lhes dá tempo de renovar o
\VUgras{isco} quando uma nova vítima se debate na ponta do fio.

%%K t08-c2-09-3 p
Com tudo isto, o \VUgras{pescador de rede} \textsuperscript{(48)}, que do seu
\VUgras{barco} \textsuperscript{(49)} lança o tarrafe no meio do \VUgras{rio}
\textsuperscript{(50)}, apanha muito mais.

%%K t08-c2-10-1 p
10. --- O outono é a estação dos bons passeios: a pé, a \VUgras{cavalo}
\textsuperscript{(52)}, de \VUgras{bicicleta} \textsuperscript{(53)}, de
\VUgras{automóvel} \textsuperscript{(54)}. Os \VUgras{caminhos} \textsuperscript{(51)}
estão limpos e aproveitam-se os últimos dias bons, porque já as
\VUgras{andorinhas} \textsuperscript{(56)} se reúnem nos telhados para voarem a
toda a asa para países mais quentes. As folhas da velha \VUgras{nogueira}
\textsuperscript{(57)} amarelecem, as \VUgras{nozes} caem e as altas
\VUgras{árvores} agrupadas em \VUgras{bosquete} \textsuperscript{(58)} sobre a
\VUgras{altura} \textsuperscript{(59)} em breve ficarão sem folhas.

%%K t08-c2-11-1 p
11. --- O \VUgras{sol} \textsuperscript{(60)} nasce mais tarde, os dias diminuem,
começa a sentir-se um frio bastante vivo de \VUgras{manhã}, e já bandos
de \VUgras{galinholas} \textsuperscript{(61)} deixam os nossos climas
inóspitos.
\VUsaut{6.0mm}
\VUfilet{22mm}
